\section{Boundary cases and consequences of realization}
\label{app:structural-consequences}

\subsection{One-edge supports with their actual masks}
\begin{lemma}\label{lem:one-edge-realization}
Let a support consist of one edge $uv$, let its directional multiplicity be $k$,
and read the actual ambient masks $M_u,M_v$. An $s$-state simple tagged cycle
exists exactly when $k\le s$ if $M_u\ne M_v$, and exactly when $2k\le s$ if
$M_u=M_v$.
\end{lemma}
\begin{proof}
Each endpoint needs $k$ distinct states. If the masks are equal, the two sets
must be disjoint because their outgoing directions are opposite. These are the
necessary inequalities. If the masks differ, use states $0,\ldots,k-1$ at both
ends and prescribe
\[
 \delta(i,M_u)=(d,i),\qquad
 \delta(i,M_v)=(-d,i+1\bmod k),
\]
where $d$ points from $u$ to $v$. If the masks agree, use disjoint sets with
$\delta(i,M)=(d,k+i)$ and $\delta(k+i,M)=(-d,i+1\bmod k)$.
These rows produce the required cycle and can be completed legally.
\end{proof}
Thus a doubled edge is possible with three states when it is the entire path,
but not when both its endpoints have the same straight ambient mask. The
multiplicity-three exception on an entire one-edge path is outside the
contour-plus-matching input of Theorem~\ref{thm:three-state-realization}.

\subsection{Two-state contours and a strict memory separation}
For a support of at least three vertices, the two-state capacity law forces
every edge multiplicity to be one. This is a necessary normal form, not an
automatic realization theorem. Fix an order of the two ports of every internal
mask $M$, and let $h_M\in\Ftwo$ be the outgoing state for its first port; the
other state is $h_M+1$. Choose endpoint states, giving at most four endpoint
profiles. At internal occurrences group equal $(M,d)$ and equate their successor
states. At endpoints group equal $(M,q)$ and do the same. The resulting equations
in the $h_M$ are scalar XOR equations, including constants.

This criterion is exact. Necessity follows because both states are used at
each internal vertex and its two outgoing directions differ. For sufficiency,
the two assigned states there are distinct, and each repeated controller key
has both its direction and successor fixed consistently by precisely the stated
groups. Complete unused rows legally. Thus one of these at most four systems is
consistent if and only if the ordinary contour is realizable with two states.
There are at most six internal masks, so this particular recognition construction
is linear in the word length for the fixed compass alphabet.

A useful strict separation is the family with directions
\[
 \E^a\N^b\W^c,\qquad a,c\ge2,\quad b\ge3.
\]
Its ordinary contour requires exactly three states. For the two-state
impossibility, let $\alpha$ be the N-state of the $\N\W$ corner and $\beta$ the
S-state of the $\Sdir\W$ corner. The shared E-row of the straight horizontal
mask forces $\alpha=\beta$, by comparing its transitions into the two corners.
The straight vertical N-row compared with the upper corner forces its state to
be $1-\beta$; the vertical S-row compared with the lower corner forces its state
to be $1-\alpha$. These are then equal, contradicting the two directions at a
straight vertical vertex. For sufficiency, the following used rows give one
three-state controller for the whole family; blank rows may be completed legally:
\begin{center}
\begin{tabular}{c@{\qquad}ccc}
\toprule
Mask & State 0 & State 1 & State 2\\
\midrule
$\E\W$ & $(\E,0)$ & $(\W,1)$ & ---\\
$\N\Sdir$ & --- & $(\Sdir,1)$ & $(\N,2)$\\
$\N\W$ & $(\N,2)$ & $(\W,1)$ & ---\\
$\Sdir\W$ & $(\Sdir,1)$ & --- & $(\W,1)$\\
$\E$ & --- & $(\E,0)$ & ---\\
\bottomrule
\end{tabular}
\end{center}
Start at the lower left endpoint in state 1. The horizontal directions use
states 0 and 1, the vertical directions states 2 and 1, and the two corners
use states $(0,1)$ and $(2,0)$. Directly following these rows gives the contour
and distinct tags at each vertex. Minimality of its eight-vertex representative
among all two-state contour failures is only a finite-check assertion.

\subsection{Proper intervals retain ambient ports}
\label{app:boundary-aware}
Suppose the prescribed support is an interval $I$ of an ambient induced path
$P$, with at least three vertices. Masks must still be read in $P$. Use the
internal profiles of Section~\ref{sec:realization} and choose injective endpoint
states. When an endpoint mask $M$ also occurs internally, require
\[
 X_i=a_M\quad\Longleftrightarrow\quad d_i=u_M.
\]
For a paired-side endpoint, its chosen state is $b(a_M)+a_Mt_i$ for a determined
constant $t_i\in\Ftwo$. Include these endpoint occurrences in the same U/B
comparisons as internal occurrences. If a mask occurs only at endpoints, group
by $(M,X_i)$, reject a group with differing outgoing directions, and otherwise
equate successors in that group. The one-edge case is Lemma~\ref{lem:one-edge-realization}.

These additions give the exact boundary-aware version of the affine criterion.
For necessity, the real controller gives the same singleton/paired partition
whenever that mask occurs internally; it therefore obeys the direction gate
and all comparisons. For sufficiency the gate checks directions at shared
endpoint/internal keys, U/B checks their successors, and the remaining endpoint
groups check all boundary-only keys. This exhausts every repeated row, while
all prescribed directions stay in $I$. Legal completion is performed using the
ambient masks. The system still has linear size and a bounded profile family.
Deleting the exterior ports instead would change the realization question.

For a fixed known controller, its internal direction partition, singleton state
and paired successor difference are already fixed. Equate each actual successor
to its known row output (affine on the two paired inputs) and impose the endpoint
rows directly. At most 36 endpoint injections remain. This is an exact linear
recognition test for a prescribed support and word in that fixed controller,
not a solution of the basin problem over all paths.

\subsection{Full-support cycles and their residual competitors}
\label{app:basins}
\begin{proposition}\label{prop:residual-cycles}
Let a three-state controller have a recurrent cycle $C$ using every vertex of a
path $P$ with $n\ge3$. Every other recurrent cycle has a one-edge support and
length two. If $R_v$ denotes the states unused by $C$ at $v$, these other cycles
are exactly the pairs
\[
 (u,q)\longmapsto(v,r)\longmapsto(u,q),\qquad
 uv\in E(P),\quad q\in R_u,\quad r\in R_v.
\]
If the vertices unsaturated by $C$ form an independent set, $C$ is the unique
cycle and every tag enters it within $3n-|C|$ steps.
\end{proposition}
\begin{proof}
Distinct cycles of a deterministic map have disjoint tags. The full-support
cycle uses at least two states at each internal vertex. A second cycle with
at least two support edges would need two further states at an internal vertex
of its support, which is also internal in $P$, an impossibility. A one-edge
competitor has an endpoint internal in $P$, where at most one state remains;
its multiplicity is therefore one. This proves the exact residual-row test.
If the unsaturated vertices are independent, no support edge can have residual
states at both ends, so no competitor exists. A trajectory outside $C$ cannot
repeat a tag before entering $C$, since repetition would give another cycle.
There are $3n-|C|$ tags outside it, giving the stated entrance bound.
\end{proof}
When all internal vertices are saturated the independent-set hypothesis holds.
In the three-state sharp-period family of Section~\ref{sec:sharp-period}, this
gives entrance within two steps for even $n$ and three for odd $n$. Once both
agents are on the unique cycle, path phase contact applies. This consequence
concerns that one explicit family of paths, not all path geometries.

The accessibility limitation is concrete already on the four-vertex path with
directions $\E\N\E$. Consider the same cycle in two controllers,
\[
 C=(0_2,1_1,2_1,3_1,2_2,1_2).
\]
Their rows in states 1 and 2 agree; state 0 differs as follows:
\begin{center}
\begin{tabular}{ccccc}
\toprule
Mask & State 1 & State 2 & State 0, good & State 0, bad\\
\midrule
$\E$ & $(\E,1)$ & $(\E,1)$ & $(\E,1)$ & $(\E,0)$\\
$\N\W$ & $(\N,1)$ & $(\W,2)$ & $(\N,1)$ & $(\W,0)$\\
$\E\Sdir$ & $(\E,1)$ & $(\Sdir,2)$ & $(\E,1)$ & $(\E,0)$\\
$\W$ & $(\W,2)$ & $(\W,2)$ & $(\W,2)$ & $(\W,0)$\\
\bottomrule
\end{tabular}
\end{center}
At every other nonempty mask choose its first port in a fixed compass order and
next state 0, completing both tables. The good controller sends every state-zero
start to $C$ in one step; the other unused tags do so too. For the bad controller,
$0_0\leftrightarrow1_0$ and $2_0\leftrightarrow3_0$ are accessible cycles, whereas
$C$ is not state-zero accessible. Starts $0_0,2_0$ alternate between spatial
pairs $(0,2)$ and $(1,3)$, always at distance two. Thus the same realized cycle
can coexist with different basin behavior.

\subsection{Why path support is essential}\label{app:branching-example}
Path support cannot be replaced by arbitrary tree support, even if the one-agent functional graph has a unique recurrent tagged cycle.  For an explicit example, take the seven-cell triod
\[
 c=(0,0),\quad n=(0,1),\quad n_2=(0,2),\quad
 s=(0,-1),\quad s_2=(0,-2),\quad e=(1,0),\quad e_2=(2,0),
\]
and the following three-state rows; unused rows may be completed arbitrarily.

\begin{center}
\small
\begin{tabular}{@{}c|ccc@{}}
\toprule
Mask & state $0$ & state $1$ & state $2$\\
\midrule
$\N\E\Sdir$ & $(\N,0)$ & $(\Sdir,1)$ & $(\E,0)$\\
$\N\Sdir$ & $(\N,2)$ & $(\Sdir,1)$ & $(\N,2)$\\
$\E\W$ & $(\E,0)$ & $(\W,0)$ & $(\E,0)$\\
$\Sdir$ & $(\Sdir,1)$ & $(\Sdir,1)$ & $(\Sdir,1)$\\
$\N$ & $(\N,0)$ & $(\N,0)$ & $(\N,0)$\\
$\W$ & $(\W,1)$ & $(\W,1)$ & $(\W,1)$\\
\bottomrule
\end{tabular}
\end{center}


These rows give the unique recurrent tagged cycle
\[
 c_0,n_0,(n_2)_2,n_1,c_1,s_1,(s_2)_1,s_0,c_2,e_0,(e_2)_0,e_1.
\]
Every off-cycle tag enters it in one step.  Nevertheless the state-zero starts $n_0$ and $s_0$, which are six phases apart, remain at Manhattan distance exactly two at all twelve phases.  Thus branching invalidates the one-dimensional sign argument and marks a real obstruction to extending the two-state method directly to three states.

